\chapter{Development of the Virtual Laboratory}

\section{Development Environment and Technologies}

The virtual laboratory was developed using a collection of modern web technologies and software development tools that support browser-based Virtual Reality (VR), real-time collaboration, and cross-platform deployment. The selected technologies were chosen to provide high rendering performance, modular software development, efficient multiplayer communication, and compatibility with WebXR-enabled devices.

Table~\ref{tab:development_technologies} summarises the principal technologies used during the development of the proposed virtual laboratory.

\begin{table}[H]
\centering
\caption{Development environment and technologies}
\label{tab:development_technologies}

\renewcommand{\arraystretch}{1.3}

\begin{tabular}{p{4cm}p{4cm}p{7cm}}
\toprule
\textbf{Category} & \textbf{Technology} & \textbf{Purpose} \\
\midrule

Programming Language &
JavaScript &
Client-side application development and interactive functionality. \\

3D Graphics Engine &
Babylon.js &
Rendering of the virtual environment, scene management, lighting, materials, animations, and WebXR integration. \\

Virtual Reality &
WebXR Device API &
Provides immersive VR support and interaction through compatible web browsers and VR headsets. \\

Backend Runtime &
Node.js &
Executes the backend server and manages multiplayer services. \\

Multiplayer Framework &
Colyseus &
Manages multiplayer rooms, state synchronisation, and real-time communication between connected users. \\

Communication Protocol &
WebSocket &
Provides persistent bidirectional communication between clients and the server. \\

Development Environment &
Visual Studio Code &
Source code development, debugging, and project management. \\

Version Control &
Git &
Source code version management and collaborative development. \\

Repository Hosting &
GitHub &
Source code hosting and project version tracking \cite{thesisRepo}. \\

Development VPS Deployment &
Nginx, PM2 &
Nginx served as the reverse proxy on the VPS used during development and evaluation; PM2 managed the Node.js backend process, providing supervision and automatic restart. \\

Self-Hosting Deployment &
Docker, Docker Compose, Caddy &
Containerised deployment configuration included in the public repository for reproducible self-hosting. Caddy provides HTTPS support and routing of frontend and backend services. \\

Target Platform &
Meta Quest 3, Desktop Browsers &
Execution of the virtual laboratory in immersive VR and desktop modes. \\

\bottomrule
\end{tabular}

\end{table}

The implementation follows a browser-based client--server architecture in which the frontend is responsible for rendering the virtual environment and handling user interaction, while the backend manages multiplayer sessions and synchronises shared virtual objects. The use of modern web technologies allows the application to operate across multiple platforms without requiring dedicated native software installation.

Two deployment configurations are relevant to this work. The virtual laboratory was developed and evaluated on a Virtual Private Server (VPS) running Nginx as a reverse proxy with PM2 managing the Node.js backend process. The public repository additionally includes a self-contained Docker Compose configuration in which Caddy provides HTTPS termination and request routing, enabling straightforward self-hosting by third parties without modifying application source code.

The complete implementation developed as part of this thesis is maintained in a public GitHub repository, enabling reproducibility, future development, and long-term maintenance of the virtual laboratory platform \cite{thesisRepo}. The repository contains the frontend application, backend server, deployment configuration, project documentation, and implementation of the interactive physics experiments.

\section{Virtual Laboratory Layout}

The virtual laboratory was designed to provide a structured and intuitive learning environment that supports both theoretical instruction and practical experimentation. Rather than representing a single room, the environment is organised into multiple interconnected spaces, each serving a specific educational or functional purpose. This organisation allows users to move naturally between different learning activities while maintaining a coherent virtual experience.

A central corridor connects all individual spaces of the environment. From this corridor, users can access the physics laboratory, the data centre, the AI laboratory, and the outdoor area without requiring additional loading screens or scene transitions. This design provides clear spatial organisation and allows users to easily understand the relationship between different functional areas of the virtual laboratory.

The physics laboratory is the primary educational space of the virtual environment. Interactive physics experiments are situated within this room; in the current implementation, it houses the Stern-Gerlach experiment, allowing users to manipulate experimental components and observe the resulting physical phenomena. The open layout provides sufficient space for multiple users to collaborate simultaneously while maintaining clear visibility of the experimental setup.

The data centre provides a second dedicated experimental space and hosts the Convex Lens experiment. Users can enter this area from the corridor and interact with the optical components of the experiment. Its separation from the main physics laboratory reflects the modular design of the environment, in which different experimental activities occupy distinct spaces that users can navigate freely.

The AI laboratory is a dedicated space for the AI-driven experiment generator feature. Rather than hosting a pre-defined experiment, it provides a space into which AI-generated experiments are instantiated at runtime. Its visual design distinguishes it from the other rooms and signals its distinct function within the virtual laboratory.

The outdoor area provides an optional environment that supports Gaussian-splat scene rendering, illustrating how the virtual laboratory can be extended beyond conventional interior spaces.

The layout was intentionally designed to support collaborative learning. Multiple users can occupy the same environment, navigate independently between spaces, and participate in shared laboratory activities. The spatial separation of different functional areas reduces visual clutter and provides a logical organisation that improves orientation within the virtual environment.

The modular organisation of the virtual laboratory facilitates future expansion. Additional laboratories or experimental environments can be integrated by extending the existing layout while preserving the overall navigation concept. Consequently, the virtual laboratory serves not only as a platform for the experiments implemented in this thesis but also as a foundation for future educational applications and collaborative learning scenarios.

The implementation of the virtual laboratory layout and its associated assets is available in the accompanying project repository \cite{thesisRepo}.


\section{User Navigation and Interaction}

An intuitive navigation and interaction mechanism is essential for providing an effective learning experience within an immersive virtual environment. The proposed virtual laboratory was designed to support both desktop users and Virtual Reality (VR) users, enabling participants to interact with the same environment regardless of the hardware platform used. The interaction model emphasises simplicity, accessibility, and natural manipulation of virtual objects while maintaining consistency across different input devices. Figure~\ref{fig:user_interaction} illustrates the overall interaction workflow within the virtual laboratory.

Upon launching the application, users enter the virtual laboratory through a standard web browser. Desktop users access the environment using conventional keyboard and mouse controls, whereas users equipped with a compatible VR headset can activate immersive mode through the integrated ``Enter VR'' functionality provided by the browser. This dual-mode approach enables both desktop and VR participants to join the same collaborative laboratory session.

Navigation within the virtual environment is designed to be straightforward and intuitive. Users are free to move throughout the virtual laboratory, allowing them to transition between the physics laboratory, data centre, AI laboratory, corridor, and outdoor area. The spatial organisation of the environment enables users to naturally explore different learning areas without interruption or scene changes, thereby providing a continuous immersive experience.

Interaction with laboratory equipment follows a direct manipulation paradigm. Interactive objects respond to user input through object selection, grabbing, movement, and release operations. In desktop mode, these interactions are performed using the keyboard and mouse, whereas in immersive mode they are executed using handheld VR controllers. Visual feedback is provided during object manipulation, allowing users to clearly identify interactive components and observe the effects of their actions.

The user interface was intentionally designed to remain simple and unobtrusive in order to preserve immersion. Only essential interface elements, such as entering VR mode and experiment-related controls, are presented to the user. Most interactions occur directly within the three-dimensional environment, allowing learners to focus on the experimental activities rather than navigating complex menus.

The virtual laboratory also supports collaborative interaction by representing every connected participant as an avatar within the shared environment. Users can observe the movements and interactions of other participants in real time, enabling collaborative experimentation, instructor demonstrations, and group discussions. Registered shared objects are synchronised across all connected clients through the Colyseus state mechanism, while certain experiment actions are communicated through discrete event messages. This ensures that every participant observes a consistent representation of the experiment.

Overall, the navigation and interaction design aims to provide an immersive, intuitive, and collaborative user experience while remaining sufficiently flexible to accommodate future experiments and additional interaction techniques. The implementation of the interaction mechanisms is available in the accompanying project repository \cite{thesisRepo}.


\begin{figure}[H]
\centering
\begin{tikzpicture}[
node distance=1.5cm,
box/.style={
draw,
rounded corners,
minimum width=4.2cm,
minimum height=0.9cm,
align=center,
fill=gray!10
},
arrow/.style={->,thick}
]

\node[box] (user) {User};

\node[box, below left=1.5cm and 2.3cm of user] (desktop) {Desktop Mode\\Keyboard \& Mouse};

\node[box, below right=1.5cm and 2.3cm of user] (vr) {VR Mode\\WebXR Controllers};

\node[box, below=3cm of user] (navigation) {Navigation};

\node[box, below left=1.5cm and 2.3cm of navigation] (interaction) {Object Interaction};

\node[box, below right=1.5cm and 2.3cm of navigation] (collaboration) {Collaborative Interaction};

\node[box, below=3cm of navigation] (lab) {Virtual Laboratory};

\draw[arrow] (user)--(desktop);
\draw[arrow] (user)--(vr);

\draw[arrow] (desktop)--(navigation);
\draw[arrow] (vr)--(navigation);

\draw[arrow] (navigation)--(interaction);
\draw[arrow] (navigation)--(collaboration);

\draw[arrow] (interaction)--(lab);
\draw[arrow] (collaboration)--(lab);

\end{tikzpicture}

\caption{User navigation and interaction workflow within the virtual laboratory.}
\label{fig:user_interaction}
\end{figure}

\section{VR Controller Interaction}

The proposed virtual laboratory was designed to provide a natural and intuitive interaction experience for users accessing the application through Virtual Reality (VR) headsets. By utilising the WebXR Device API together with the interaction capabilities provided by Babylon.js, users are able to manipulate virtual laboratory equipment using handheld VR controllers. This interaction model closely resembles real-world object manipulation and contributes to a more immersive learning experience. Figure~\ref{fig:vr_controller_interaction} illustrates the interaction workflow implemented for VR users.

When a user enters immersive mode, the virtual representations of the handheld controllers are automatically initialised within the three-dimensional environment. The position and orientation of each controller are continuously updated using six degrees of freedom (6DoF) tracking provided by the VR headset. This enables users to naturally point, reach, and interact with virtual objects while moving within the laboratory.

Object interaction is implemented using Babylon.js's \texttt{SixDofDragBehavior}, which provides 6DoF dragging for both VR controller and desktop pointer input. Users can select interactive laboratory components by pointing at them and activating the grab action; the selected object then follows the controller position and orientation in six degrees of freedom. Once grabbed, it can be repositioned and released within the virtual environment. These interactions are used to manipulate experimental equipment such as optical components during physics experiments.

To provide immediate feedback during interaction, visual cues are displayed whenever an interactive object is hovered over or actively held. This feedback assists users in identifying objects that can be interacted with and confirms successful execution of interaction events. Such visual responses improve usability while reducing ambiguity during immersive experimentation.

The VR controller interaction system is integrated with the multiplayer synchronisation mechanism of the virtual laboratory through an explicit object-ownership model. When a user initiates a grab, a grab request is sent to the server; the server grants exclusive ownership by setting the \texttt{ownerId} field on the corresponding shared object, and only the owning client may subsequently submit transform updates. Position and orientation changes are transmitted to the server at approximately 20 Hz while the object is held, and the server propagates the updated state to all other connected participants. On release, ownership is cleared so that other users may subsequently interact with the object.

The interaction model was designed to remain extensible for future laboratory experiments. Additional interaction techniques, experimental devices, and educational activities can be incorporated without requiring fundamental changes to the underlying controller framework. The implementation of the VR controller interaction system is available in the accompanying project repository \cite{thesisRepo}.

\begin{figure}[H]
\centering
\begin{tikzpicture}[
node distance=1.5cm,
box/.style={
draw,
rounded corners,
minimum width=4.2cm,
minimum height=0.9cm,
align=center,
fill=gray!10
},
arrow/.style={->,thick}
]

\node[box] (controller) {VR Controller};

\node[box, below=of controller] (tracking) {6DoF Tracking};

\node[box, below=of tracking] (selection) {Object Selection};

\node[box, below left=1.5cm and 2.2cm of selection] (grab) {Grab / Move};

\node[box, below right=1.5cm and 2.2cm of selection] (release) {Release};

\node[box, below=3cm of selection] (sync) {Object Synchronisation};

\node[box, below=of sync] (scene) {Virtual Laboratory};

\draw[arrow] (controller)--(tracking);
\draw[arrow] (tracking)--(selection);
\draw[arrow] (selection)--(grab);
\draw[arrow] (selection)--(release);
\draw[arrow] (grab)--(sync);
\draw[arrow] (release)--(sync);
\draw[arrow] (sync)--(scene);

\end{tikzpicture}

\caption{VR controller interaction workflow within the virtual laboratory.}
\label{fig:vr_controller_interaction}
\end{figure}

\section{Browser and VR Device Support}

A fundamental design objective of the proposed virtual laboratory is to provide broad accessibility by supporting both conventional desktop browsers and immersive Virtual Reality (VR) devices. Rather than developing separate applications for desktop and VR platforms, a single browser-based application was implemented using modern web technologies and the WebXR Device API. This unified approach allows users to access the same virtual laboratory regardless of the hardware platform being used.

Desktop users access the virtual laboratory through a modern web browser using conventional keyboard and mouse controls. This mode provides complete access to the virtual environment, allowing users to explore laboratory spaces, participate in collaborative sessions, and interact with experimental equipment without requiring dedicated VR hardware. Supporting desktop access improves the accessibility of the system, particularly in educational institutions where VR headsets may not be available for every student.

For immersive interaction, the virtual laboratory supports WebXR-compatible Virtual Reality headsets. In this implementation, the Meta Quest 3 headset was used as the primary development and testing platform. When a compatible VR device is detected, users can enter immersive mode directly from the browser without installing additional software. Once the VR session begins, the application automatically switches to stereoscopic rendering and enables interaction using handheld VR controllers.

Both desktop and VR users access the same application, virtual environment, and multiplayer session. The interaction model is adapted according to the available input devices while maintaining identical laboratory content and collaborative functionality. Consequently, users employing different hardware platforms can participate simultaneously within the same virtual laboratory session, observe each other's activities, and interact with shared experimental equipment.

The browser-based deployment model also simplifies software distribution and maintenance. Since the application is hosted on a web server, users always access the latest available version without performing manual software installation or updates. This approach reduces deployment complexity while ensuring consistent functionality across supported devices.

Although the virtual laboratory was primarily developed and evaluated using the Meta Quest 3 headset, the implementation follows the WebXR standard and is therefore expected to operate on other compatible VR devices supporting the WebXR Device API. Similarly, desktop compatibility is achieved through standards-compliant web browsers that support the required web technologies. The implementation supporting browser and VR device compatibility is available in the accompanying project repository \cite{thesisRepo} and utilises the WebXR capabilities provided by Babylon.js \cite{babylonWebXR}.

\section{Performance Considerations}

Performance is a critical aspect of browser-based Virtual Reality (VR) applications, as smooth rendering and responsive interaction are essential for providing an immersive user experience. During the development of the proposed virtual laboratory, several design decisions were made with performance in mind, covering rendering architecture, scene organisation, and network communication.

A browser-based architecture was adopted to simplify deployment while utilising hardware-accelerated rendering through WebGL. Babylon.js provides efficient scene management, rendering optimisation, and resource handling, enabling the application to render complex three-dimensional environments within modern web browsers. The use of WebXR further allows immersive rendering without requiring dedicated native applications.

The virtual laboratory environment was designed using a modular scene structure in which static assets, interactive objects, and educational content are organised independently. A zone-based room manager enables or disables entire room sub-trees depending on the player's current position, so only the rooms a user can currently see are rendered. Static environmental assets remain unchanged during execution, whereas only interactive objects and user-controlled entities require continuous updates.

Multiplayer communication uses different synchronisation mechanisms depending on the nature of the data. Player pose updates, including position, orientation, and controller laser state, are sent by each client at 20 Hz. Manipulable shared objects such as the Convex Lens components are registered in the Colyseus shared state and updated at approximately 20 Hz while held by a user, so all participants observe consistent object positions. Discrete experiment actions, such as firing or resetting the Stern-Gerlach apparatus, are communicated through targeted event messages broadcast to all other participants. Static resources such as models, textures, and environment geometry are loaded locally by each client during initialisation, avoiding redundant server transmission.

Rendering performance is also supported by separating client-side rendering from server-side synchronisation. The client is responsible for graphics rendering, user interaction, and local visual effects, whereas the server maintains the authoritative state of the collaborative environment. This separation distributes computational workload across the system and enables each component to perform specialised tasks independently.

The browser-based deployment model further contributes to maintainability by allowing users to access the latest application version directly from the web server without manual software installation or updates. Since application assets are loaded through standard web technologies, future optimisations such as asset compression, model optimisation, and texture streaming can be incorporated without modifying the overall system architecture.

Although the primary objective of this thesis was the implementation of a functional collaborative virtual laboratory rather than extensive performance optimisation, the adopted architecture provides a foundation for future improvements. Additional optimisation techniques, including model simplification, level-of-detail (LOD) rendering, and advanced network synchronisation strategies, can be incorporated as the virtual laboratory evolves.

The implementation of the performance-related design decisions described in this section is available in the accompanying project repository \cite{thesisRepo}.
